\section{Preliminaries and the centered homotopy operator}

We begin by fixing notation and recalling the midpoint-centered radial
homotopy operator in the form needed below.

The ambient dimension in the homotopy sections is denoted by $n$, while the
degree of the active coordinate-monomial form is denoted by $k$. In Section~4,
where the top form on a rectangular box is the main object, the same letter $k$
becomes the dimension of that active box. This convention matches the sliced
applications: inactive coordinates may be present in the surrounding ambient
cube, but the sharp boundary-readout theorem is always stated in the active
top-form dimension.

Fix $n\ge 1$ and write
\[
I^n:=[0,1]^n,
\qquad
c:=\left(\frac12,\dots,\frac12\right).
\]
For $1\le k\le n$, let $\Omega^k_{\mathrm{sm}}(I^n)$ denote the smooth
$k$-forms on a neighborhood of $I^n$, restricted to the cube. If
$R\subset\RR^n$ is any axis-aligned box, we use $\Omega^k_{\mathrm{sm}}(R)$ for
the analogous space of smooth $k$-forms on a neighborhood of $R$, restricted
to $R$. If $F$ is a codimension-one face of such a box and
$\iota_F:F\hookrightarrow R$ denotes the inclusion, we write
\[
\Tr_F(\eta):=\iota_F^\ast\eta
\]
for the face trace of a differential form $\eta$ on $R$. If
\[
\omega=\sum_{\abs{I}=k} a_I(x)\,\dd x_I,
\]
we set
\[
\norm{\omega}_{\coeff,\infty}:=
\max_{\abs{I}=k}\sup_{x\in I^n}\abs{a_I(x)}.
\]
All $L^\infty$ estimates below are taken with respect to this coordinate norm.
For a form evaluated at a fixed point $y$, we write
\[
\abs{\omega(y)}_{\coeff,\infty}:=
\max_{\abs{I}=k}\abs{a_I(y)}.
\]
If
\[
\beta=\sum_{\abs{J}=m} b_J(x)\,\dd x_J,
\]
we also write $\coeff_J(\beta):=b_J$ for the coefficient of $\dd x_J$.

\begin{remark}
Throughout, $\dd$ denotes the exterior derivative on smooth forms, whereas
$\cob$ denotes the cubical coboundary used in Section~5.
\end{remark}

\begin{definition}[Coordinate-monomial form]
\label{def:decomposable-form}
A smooth $k$-form $\omega\in\Omega^k_{\mathrm{sm}}(I^n)$ is
\emph{coordinate-monomial} if either $\omega=0$, or there are an ordered
coordinate set
\[
I=\{i_1<\cdots<i_k\}\subset\{1,\dots,n\}
\]
and a smooth coefficient $f$ on a neighborhood of $I^n$ such that
\[
\omega=f(x)\,\dd x_I
:=
f(x)\,\dd x_{i_1}\wedge\cdots\wedge \dd x_{i_k}.
\]
\end{definition}

\begin{remark}
This requirement is stronger than decomposability in
$\Lambda^kT^\ast I^n$. For instance,
$(\dd x_1+\dd x_2)\wedge \dd x_3$ is decomposable in the exterior algebra, but
it is not coordinate-monomial.
\end{remark}

Let $X(x):=x-c$ be the radial vector field centered at $c$, and let
\[
H:[0,1]\times I^n\to I^n,
\qquad
H(t,x):=c+t(x-c).
\]
For each $t\in[0,1]$ we write $H_t(x):=H(t,x)$.

\begin{definition}[Radial homotopy operator]
\label{def:homotopy-operator}
For $1\le k\le n$ define
\[
K_k:\Omega^k_{\mathrm{sm}}(I^n)\longrightarrow
\Omega^{k-1}_{\mathrm{sm}}(I^n)
\]
by
\begin{equation}
\label{eq:K-definition}
(K_k\omega)(x):=
\int_0^1 t^{k-1}\,
\bigl(\iota_{x-c}\omega\bigr)\bigl(c+t(x-c)\bigr)\,\dd t .
\end{equation}
Equivalently, if
\[
\omega=\sum_{\abs{I}=k}a_I(y)\,\dd y_I,
\qquad I=\{i_1<\cdots<i_k\},
\]
then
\begin{equation}
\label{eq:K-coefficient-formula}
(K_k\omega)(x)=
\int_0^1 t^{k-1}
\sum_{\abs{I}=k}\sum_{r=1}^k
(-1)^{r-1}(x_{i_r}-c_{i_r})\,
a_I(c+t(x-c))\,\dd x_{I\setminus\{i_r\}}\,\dd t .
\end{equation}
We also set $K_{n+1}:=0$.
\end{definition}

This is the standard radial homotopy formula centered at the midpoint of the
cube; see \cite[Sec.~3]{BottTu1982} and
\cite{Dupont1976,DupontBook1978,Albert2020}. In particular, for
$\omega_0=\dd x_1\wedge\cdots\wedge \dd x_k$ one obtains
\begin{equation}
\label{eq:explicit-minimizer}
K_k\omega_0
=
\frac{1}{k}\sum_{j=1}^k (-1)^{j-1}\left(x_j-\frac12\right)
\dd x_1\wedge\cdots\wedge\widehat{\dd x_j}\wedge\cdots\wedge \dd x_k.
\end{equation}

\begin{lemma}[Contraction estimate]
\label{lem:contraction-estimate}
Let
\[
\omega=f(y)\,\dd x_{i_1}\wedge\cdots\wedge\dd x_{i_k}
\]
be coordinate-monomial. Then for all $x,y\in I^n$,
\[
\abs{\bigl(\iota_{x-c}\omega\bigr)(y)}_{\coeff,\infty}
\le \frac12\,\norm{\omega}_{\coeff,\infty}.
\]
Consequently,
\[
\norm{\iota_X\omega}_{\coeff,\infty}
\le \frac12\,\norm{\omega}_{\coeff,\infty}.
\]
\end{lemma}

\begin{proof}
Writing $I=\{i_1<\cdots<i_k\}$, contraction gives
\[
\bigl(\iota_{x-c}\omega\bigr)(y)
=
\sum_{r=1}^k
(-1)^{r-1}(x_{i_r}-c_{i_r})f(y)\,
\dd x_{I\setminus\{i_r\}} .
\]
Every coefficient therefore has absolute value at most
$\frac12\abs{f(y)}$, since $c_i=1/2$ and $0\le x_i\le 1$. Taking the maximum
over the output coefficients and then the supremum over $y$ proves the first
estimate. The second estimate is the special case $y=x$ followed by taking
the supremum over $x\in I^n$.
\end{proof}

\begin{proposition}[Incidence criterion]
\label{prop:incidence-criterion}
Let
\[
\omega=\sum_{\abs{I}=k}a_I(x)\,\dd x_I\in\Omega^k_{\mathrm{sm}}(I^n).
\]
For each ordered $(k-1)$-index $J$ set
\[
m_J(\omega):=\#\{i\notin J:\ a_{J\cup\{i\}}\not\equiv 0\},
\qquad
m_\ast(\omega):=\max_{\abs{J}=k-1}m_J(\omega),
\]
where $J\cup\{i\}$ is always written in increasing order. Then
\begin{equation}
\label{eq:incidence-contraction-bound}
\norm{\iota_X\omega}_{\coeff,\infty}
\le
\frac{m_\ast(\omega)}2\norm{\omega}_{\coeff,\infty}.
\end{equation}
In particular, let $\mathcal M$ be a family of $k$-indices and let
$\Omega^k_{\mathcal M}(I^n)$ denote the subspace of forms
\[
\omega=\sum_{I\in\mathcal M}a_I(x)\,\dd x_I .
\]
The universal estimate
\[
\norm{\iota_X\omega}_{\coeff,\infty}
\le
\frac12\norm{\omega}_{\coeff,\infty}
\qquad
(\omega\in\Omega^k_{\mathcal M}(I^n))
\]
holds if and only if every $(k-1)$-index $J$ is contained in at most one member
of $\mathcal M$. This is only the contraction-level criterion; the corresponding
radial-homotopy obstruction with the sharp factor $1/(2k)$ is stated separately
in Proposition~\ref{prop:K-incidence-criterion}.
\end{proposition}

\begin{proof}
Fix an ordered $(k-1)$-index $J$. The coefficient of $\dd x_J$ in
$\iota_X\omega$ is
\[
\sum_{\substack{i\notin J\\ I=J\cup\{i\}}}
(-1)^{r(i,I)-1}(x_i-\tfrac12)a_I(x),
\]
where $r(i,I)$ is the position of $i$ in the ordered set $I$. Each summand has
absolute value at most $\frac12\norm{\omega}_{\coeff,\infty}$ on $I^n$, and at
most $m_J(\omega)$ summands are nonzero. Taking the maximum over $J$ proves
\eqref{eq:incidence-contraction-bound}.

If each $(k-1)$-index $J$ is contained in at most one member of $\mathcal M$,
then every output coefficient of $\iota_X\omega$ has at most one contributing
input coefficient. The same computation therefore gives the bound with the
factor $1/2$. After integration in the homotopy parameter this is exactly the
one-incidence mechanism used in Proposition~\ref{prop:K-incidence-criterion}.

Conversely, suppose that some $(k-1)$-index $J$ is contained in two distinct
members $I_1=J\cup\{p\}$ and $I_2=J\cup\{q\}$ of $\mathcal M$. Define a constant
form
\[
\omega:=(-1)^{r(p,I_1)-1}\dd x_{I_1}
+(-1)^{r(q,I_2)-1}\dd x_{I_2}.
\]
Then $\norm{\omega}_{\coeff,\infty}=1$. At the point whose $p$- and
$q$-coordinates are equal to $1$, the coefficient of $\dd x_J$ in
$\iota_X\omega$ is
\[
(1-\tfrac12)+(1-\tfrac12)=1.
\]
Hence $\norm{\iota_X\omega}_{\coeff,\infty}\ge1>1/2$, so the universal
one-half estimate fails. This proves the criterion.

\end{proof}

\begin{proposition}[Incidence criterion for the radial homotopy]
\label{prop:K-incidence-criterion}
Let $\mathcal M$ be a family of $k$-indices and set
\[
\Omega^k_{\mathcal M}(I^n):=
\left\{\sum_{I\in\mathcal M}a_I(x)\,\dd x_I\right\}.
\]
Define
\[
m^*(\mathcal M):=
\max_{\abs{J}=k-1}\#\{I\in\mathcal M:\ J\subset I\}.
\]
Then every $\omega\in\Omega^k_{\mathcal M}(I^n)$ satisfies
\[
\norm{K_k\omega}_{\coeff,\infty}
\le
\frac{m^*(\mathcal M)}{2k}\norm{\omega}_{\coeff,\infty}.
\]
Moreover, the universal estimate
\[
\norm{K_k\omega}_{\coeff,\infty}
\le
\frac{1}{2k}\norm{\omega}_{\coeff,\infty}
\qquad
(\omega\in\Omega^k_{\mathcal M}(I^n))
\]
holds if and only if every $(k-1)$-index is contained in at most one member of
$\mathcal M$.
\end{proposition}

\begin{proof}
Fix an ordered $(k-1)$-index $J$. By the coefficient formula
\eqref{eq:K-coefficient-formula}, the coefficient of $\dd x_J$ in $K_k\omega$
is
\[
\int_0^1t^{k-1}
\sum_{\substack{i\notin J\\ J\cup\{i\}\in\mathcal M}}
\pm (x_i-\tfrac12)a_{J\cup\{i\}}(c+t(x-c))\,\dd t,
\]
where $J\cup\{i\}$ is written in increasing order and the sign is the usual
incidence sign. There are at most $m^*(\mathcal M)$ summands, and each has
absolute value at most $\frac12\norm{\omega}_{\coeff,\infty}$. Since
$\int_0^1t^{k-1}\,\dd t=1/k$, this gives the stated
$m^*(\mathcal M)/(2k)$ bound.

If each $(k-1)$-index lies in at most one member of $\mathcal M$, then
$m^*(\mathcal M)\le1$ and the preceding estimate gives the universal
$1/(2k)$ bound. Conversely, suppose that some ordered $(k-1)$-index $J$ is
contained in two members $I_1=J\cup\{p\}$ and $I_2=J\cup\{q\}$ of
$\mathcal M$. Choose the constant form
\[
\omega:=(-1)^{r(p,I_1)-1}\dd x_{I_1}
+(-1)^{r(q,I_2)-1}\dd x_{I_2},
\]
where $r(i,I)$ is the position of $i$ in the ordered set $I$. Then
$\norm{\omega}_{\coeff,\infty}=1$. At any point $x$ with $x_p=x_q=1$, the
$\dd x_J$ coefficient of $K_k\omega$ equals
\[
\int_0^1t^{k-1}\bigl((1-\tfrac12)+(1-\tfrac12)\bigr)\,\dd t
=\frac1k,
\]
which is larger than $1/(2k)$. Hence the universal $1/(2k)$ estimate fails.
\end{proof}

\begin{remark}
\label{rem:non-coord-monomial}
The restriction to coordinate-monomial forms is essential for the coefficient
norm. On $I^3$ consider
\[
\omega:=\dd x_1\wedge \dd x_3+\dd x_2\wedge \dd x_3.
\]
Then $\norm{\omega}_{\coeff,\infty}=1$, whereas
\[
\iota_X\omega
=
\bigl(x_1-\tfrac12+x_2-\tfrac12\bigr)\dd x_3
-\bigl(x_3-\tfrac12\bigr)(\dd x_1+\dd x_2),
\]
so $\norm{\iota_X\omega}_{\coeff,\infty}=1$. A direct computation gives
\[
K_2\omega
=
\frac{x_1+x_2-1}{2}\,\dd x_3
-\frac{x_3-\tfrac12}{2}\,(\dd x_1+\dd x_2),
\]
hence $\norm{K_2\omega}_{\coeff,\infty}=\frac12>\frac14$. Thus the estimate
quoted in Section~3 is genuinely tied to the coordinate-monomial setting.
\end{remark}

